% ========================================================
% BÀI TẬP TOÁN 5 - KẾT NỐI TRI THỨC VỚI CUỘC SỐNG
% Bài 14: Luyện tập chung
% Định dạng: BTTeX Part III - Trắc nghiệm điền số ngắn
% ========================================================

% Câu 1
\begin{ex}
  Trong số thập phân $48{,}205$, chữ số $0$ ở hàng nào? Điền số tương ứng: 1 (hàng phần mười), 2 (hàng phần trăm), 3 (hàng phần nghìn).
  \shortans{2}
  \loigiai{
    Chữ số 0 đứng ở vị trí thứ hai sau dấu phẩy, thuộc hàng phần trăm (mã số 2).
  }
\end{ex}

% Câu 2
\begin{ex}
  Viết phân số $\frac{3}{5}$ thành số thập phân ta được $0{,}6$. Lấy phần nguyên cộng chữ số hàng phần mười được bao nhiêu?
  \shortans{6}
  \loigiai{
    Phần nguyên là 0, chữ số hàng phần mười là 6. Tổng là 6.
  }
\end{ex}

% Câu 3
\begin{ex}
  Tìm số tự nhiên lớn nhất bé hơn $15{,}78$. Số đó là bao nhiêu?
  \shortans{15}
  \loigiai{
    Số tự nhiên lớn nhất bé hơn $15{,}78$ là 15.
  }
\end{ex}

% Câu 4
\begin{ex}
  Tìm số tự nhiên bé nhất lớn hơn $29{,}05$. Số đó là bao nhiêu?
  \shortans{30}
  \loigiai{
    Số tự nhiên bé nhất lớn hơn $29{,}05$ là 30.
  }
\end{ex}

% Câu 5
\begin{ex}
  Có bao nhiêu số tự nhiên nằm giữa $8{,}9$ và $14{,}1$?
  \shortans{6}
  \loigiai{
    Các số tự nhiên là $9, 10, 11, 12, 13, 14$. Có tất cả $14 - 9 + 1 = 6$ số.
  }
\end{ex}

% Câu 6
\begin{ex}
  Làm tròn số $47{,}85$ đến hàng đơn vị ta được số nào?
  \shortans{48}
  \loigiai{
    Chữ số hàng phần mười là 8 nên làm tròn lên thành 48.
  }
\end{ex}

% Câu 7
\begin{ex}
  Làm tròn số $63{,}42$ đến hàng đơn vị ta được số nào?
  \shortans{63}
  \loigiai{
    Chữ số hàng phần mười là 4 nên làm tròn xuống thành 63.
  }
\end{ex}

% Câu 8
\begin{ex}
  Đổi $5\text{ m } 8\text{ cm}$ thành mét dưới dạng số thập phân là $5{,}08\text{ m}$. Chữ số hàng phần trăm là bao nhiêu?
  \shortans{8}
  \loigiai{
    $5\text{ m } 8\text{ cm} = 5{,}08\text{ m}$. Chữ số hàng phần trăm là 8.
  }
\end{ex}

% Câu 9
\begin{ex}
  Đổi $4\text{ tấn } 5\text{ kg}$ thành tấn dưới dạng số thập phân là $4{,}005\text{ tấn}$. Chữ số hàng phần nghìn là bao nhiêu?
  \shortans{5}
  \loigiai{
    $4\text{ tấn } 5\text{ kg} = 4{,}005\text{ tấn}$. Chữ số hàng phần nghìn là 5.
  }
\end{ex}

% Câu 10
\begin{ex}
  Tìm chữ số $x$ để $23{,}x7 > 23{,}87$. Chữ số $x$ là bao nhiêu?
  \shortans{9}
  \loigiai{
    Để $23{,}x7 > 23{,}87$ thì $x > 8$. Chữ số duy nhất thỏa mãn là 9.
  }
\end{ex}

% Câu 11
\begin{ex}
  Sắp xếp bốn số $4{,}5; 4{,}05; 4{,}55; 4{,}055$ theo thứ tự từ lớn đến bé. Số lớn nhất có chữ số hàng phần trăm là bao nhiêu?
  \shortans{5}
  \loigiai{
    Số lớn nhất là $4{,}55$. Chữ số hàng phần trăm là 5.
  }
\end{ex}

% Câu 12
\begin{ex}
  Đổi $3\text{ kg } 750\text{ g}$ thành ki-lô-gam dưới dạng số thập phân ta được $3{,}75\text{ kg}$. Tích hai chữ số ở phần thập phân là bao nhiêu?
  \shortans{35}
  \loigiai{
    $3{,}75$ có hai chữ số phần thập phân là 7 và 5. Tích là $7 \times 5 = 35$.
  }
\end{ex}

% Câu 13
\begin{ex}
  Đổi $2\text{ m}^2\text{ } 45\text{ dm}^2$ thành mét vuông dưới dạng số thập phân ta được $2{,}45\text{ m}^2$. Chữ số hàng phần mười là bao nhiêu?
  \shortans{4}
  \loigiai{
    $2\text{ m}^2\text{ } 45\text{ dm}^2 = 2{,}45\text{ m}^2$. Chữ số hàng phần mười là 4.
  }
\end{ex}

% Câu 14
\begin{ex}
  Một hình vuông có cạnh dài $2{,}5\text{ m}$. Chu vi hình vuông đó là bao nhiêu mét?
  \shortans{10}
  \loigiai{
    Chu vi là $2{,}5 \times 4 = 10\text{ m}$.
  }
\end{ex}

% Câu 15
\begin{ex}
  Một hình chữ nhật có chiều dài $4{,}5\text{ m}$, chiều rộng $2\text{ m}$. Diện tích hình chữ nhật đó là bao nhiêu mét vuông?
  \shortans{9}
  \loigiai{
    Diện tích là $4{,}5 \times 2 = 9\text{ m}^2$.
  }
\end{ex}

% Câu 16
\begin{ex}
  Chuyển phân số $\frac{7}{25}$ thành số thập phân ta được $0{,}28$. Tổng các chữ số của số thập phân đó là bao nhiêu?
  \shortans{10}
  \loigiai{
    $\frac{7}{25} = \frac{28}{100} = 0{,}28$. Tổng các chữ số là $0 + 2 + 8 = 10$.
  }
\end{ex}

% Câu 17
\begin{ex}
  Chuyển phân số $\frac{3}{8}$ thành số thập phân ta được $0{,}375$. Chữ số hàng phần nghìn là bao nhiêu?
  \shortans{5}
  \loigiai{
    $\frac{3}{8} = \frac{375}{1000} = 0{,}375$. Chữ số hàng phần nghìn là 5.
  }
\end{ex}

% Câu 18
\begin{ex}
  Tìm chữ số $y$ sao cho $5{,}6y8$ khi làm tròn đến hàng phần trăm được số $5{,}65$. Chữ số $y$ là bao nhiêu?
  \shortans{4}
  \loigiai{
    Vì chữ số hàng phần nghìn là 8 ($8 \ge 5$) nên được làm tròn lên. Để sau khi làm tròn được chữ số 5 thì trước đó $y = 5 - 1 = 4$.
  }
\end{ex}

% Câu 19
\begin{ex}
  Có bao nhiêu số tự nhiên $a$ thỏa mãn $4{,}2 \times 2 < a < 15{,}3$?
  \shortans{7}
  \loigiai{
    $4{,}2 \times 2 = 8{,}4$. Ta có $8{,}4 < a < 15{,}3$. Các số tự nhiên là $9, 10, 11, 12, 13, 14, 15$. Có 7 số.
  }
\end{ex}

% Câu 20
\begin{ex}
  Một bao xi măng cân nặng $49{,}8\text{ kg}$. Làm tròn khối lượng bao xi măng đến hàng đơn vị ta được bao nhiêu ki-lô-gam?
  \shortans{50}
  \loigiai{
    Chữ số hàng phần mười là 8 nên làm tròn lên thành $50\text{ kg}$.
  }
\end{ex}

% Câu 21
\begin{ex}
  Viết số thập phân bé nhất có ba chữ số khác nhau ở phần thập phân từ các chữ số $0; 3; 6; 9$ với phần nguyên là $0$. Chữ số hàng phần mười là bao nhiêu?
  \shortans{3}
  \loigiai{
    Số bé nhất là $0{,}369$. Chữ số hàng phần mười là 3.
  }
\end{ex}

% Câu 22
\begin{ex}
  Viết số thập phân lớn nhất có hai chữ số ở phần thập phân từ các chữ số $1; 4; 7$ với phần nguyên có một chữ số. Số đó là $7{,}41$. Chữ số hàng phần trăm là bao nhiêu?
  \shortans{1}
  \loigiai{
    Chữ số hàng phần trăm là 1.
  }
\end{ex}

% Câu 23
\begin{ex}
  Một người chạy quãng đường $100\text{ m}$ hết $12{,}8\text{ giây}$. Làm tròn thời gian chạy đến hàng đơn vị ta được bao nhiêu giây?
  \shortans{13}
  \loigiai{
    Chữ số hàng phần mười là 8 nên làm tròn lên thành 13 giây.
  }
\end{ex}

% Câu 24
\begin{ex}
  Đổi $0{,}6\text{ giờ}$ thành bao nhiêu phút?
  \shortans{36}
  \loigiai{
    $0{,}6\text{ giờ} = 0{,}6 \times 60 = 36\text{ phút}$.
  }
\end{ex}

% Câu 25
\begin{ex}
  Đổi $0{,}25\text{ ngày}$ thành bao nhiêu giờ?
  \shortans{6}
  \loigiai{
    $0{,}25\text{ ngày} = 0{,}25 \times 24 = 6\text{ giờ}$.
  }
\end{ex}
